\section{Conclusions}

Practical Byzantine Fault Tolerance (pBFT) successfully bridged the gap between theoretical models of Byzantine fault tolerance and real-world distributed systems. By optimizing cryptographic operations---relying primarily on efficient symmetric Message Authentication Codes (MACs) and reserving expensive public-key cryptography only for critical view-change procedures---pBFT demonstrated that resilience against arbitrary, malicious node behavior could be achieved without crippling system performance. The empirical results from the Byzantine-Fault-tolerant File System (BFS) proved that pBFT incurs merely a marginal overhead compared to standard, unreplicated systems.

Furthermore, mechanisms such as garbage collection via stable checkpoints and proactive recovery through secure hardware coprocessors address the limitations of unbounded log growth and long-term security, ensuring the protocol remains robust over extended lifespans. 

While initially considered an academic exercise due to the industry's reliance on perimeter-based security, the rise of the Zero Trust security paradigm and decentralized networks has catalyzed a renaissance for pBFT. Its ability to provide immediate transaction finality, strict linearizability, and high throughput without the massive energy consumption associated with Proof-of-Work mechanisms has established it as the foundational consensus architecture for modern permissioned blockchains, such as Hyperledger, and federated financial networks. Ultimately, pBFT remains a seminal work, proving that achieving high-performance consensus in untrusted environments is both mathematically rigorous and practically viable.